% !TeX root = ../教材题库_全部条目_转写工作稿.tex
% Source: XJ-B1-C04-EX-P108-Y0177
% 第4章 / 4.2 指数函数 / 教材 p.108 / PDF 第115页
% OCR未校对：公式与文字均可能有误；未完成数学TeX转写；禁止直接发给学生
% 原题及相邻原解截图（核对用）：../题目原图/XJ-B1-C04-EX-P108-Y0177-01.webp
% 本文件只在转写完成后作为题干使用；OCR 条目尚不是可靠题干。
\begin{bankitem}{4.2 例题 例题1}{第4章 / 4.2 指数函数 / 教材 p.108 / PDF 第115页}{XJ-B1-C04-EX-P108-Y0177}{OCR未校对：公式与文字均可能有误；未完成数学TeX转写；禁止直接发给学生}
\begin{ocrdraft}
[教材p.108 / PDF第115页]
\par
00某公司为大力提升创新水平，计划逐年加大研发资金投人．若该公司当
\par
年投人的研发资金为3500万元，预计下一年（将为计划执行的第一年）投人的研发
\par
资金为4200万元；如果假定增速不变，该公司第x年投人的研发资金用函数
\par
G(x)=C．来近似地表示，写出此函数的解析式，并以此估计计划执行的第8年
\par
该公司投人研发资金为多少万元（1.284．3).
\par
解按假设条件和数据，有
\par
G(O)=C\ensuremath{\cdot}=3500，
\par
G(D=C\ensuremath{\cdot}是=4200．
\par
解得
\par
C==3500
\par
4200
\par
4200
\par
6
\par
于是
\par
C
\par
3500
\par
5．
\par
于是，该函数解析式为G(x）一3500\ensuremath{\times}
\par
因此，可估计出该公司第8年投人的研发资金为
\par
G(8）=C．流=3500\ensuremath{\times}
\par
15050（万元).
\end{ocrdraft}
\end{bankitem}
